\section{Рабочий проект}

В процессе разработки были выделены следующие модули системы: \texttt{common}, \texttt{macro}, \texttt{compiler}, \texttt{optimizer}, \texttt{generator}, \texttt{assembler}.


\subsection{common.lsp}

\texttt{common.lsp} -- вспомогательный модуль, содержащий функции, являющиеся общими для нескольких других модулей.

Описание функций модуля \texttt{common.lsp} представлено в таблице \ref{common:table}.

\renewcommand{\arraystretch}{0.8} % уменьшение расстояний до сетки таблицы
\begin{xltabular}{\textwidth}{|p{2.5cm}|p{2.5cm}|X|}
	\caption{Описание функций, используемых в модуле \texttt{common.lsp}\label{common:table}}\\
	\hline \centrow \setlength{\baselineskip}{0.7\baselineskip} Название функции & \centrow Аргументы функции & \centrow Описание функции \\
	\hline \centrow 1 & \centrow 2 & \centrow 3 \\ \hline
	\endfirsthead
	\caption*{Продолжение таблицы \ref{common:table}}\\
	\hline \centrow 1 & \centrow 2 & \centrow 3 \\ \hline
	\finishhead
	comp-err & msg, \&rest other & Вызывает исключение \texttt{comp-err} с сообщением \texttt{msg} и опциональной дополнительной информацией об ошибке \texttt{other}. \\ \hline
	is-nary & args & Проверка на присутствие \texttt{\&rest} аргумента в \texttt{args}. \\ \hline
	mk/add-func & name list \&rest other & Макрос для генерации функции \texttt{name}, добавляющей информацию компиляции функций в список \texttt{list} с дополнительной информацией \texttt{other}. \\ \hline
	num-fix-args & list num & Определяет число фиксированных аргументов в списке \texttt{list} с начальным значением \texttt{num}. \\ \hline
	remove-rest & list & Удаляет символ \texttt{\&rest} из списка аргументов \texttt{list}. \\ \hline
	make-nary-args & count args & Помещает необязательные аргументы из списка \texttt{args}, в котором \texttt{count} обязательных аргументов, в ещё один список. \\ \hline
	search-symbol & list name & Поиск информации о функции или примитиве с названием \texttt{name} в списке \texttt{list}. \\ \hline
	correct-lambda & f & Проверка на то, что S-выражение \texttt{f} является корректным лямбда-выражением. \\ \hline
\end{xltabular}
\renewcommand{\arraystretch}{1.0} % восстановление сетки


\subsection{macro.lsp}

\texttt{macro.lsp} -- модуль, отвечающий за регистрацию новых и раскрытие существующих макросов. Используется модулем \texttt{compiler.lsp} и преобразовывает S-выражения, осуществляя макроподстановку.

Описание функций модуля \texttt{macro.lsp} представлено в таблице \ref{macro:table}.

\renewcommand{\arraystretch}{0.8} % уменьшение расстояний до сетки таблицы
\begin{xltabular}{\textwidth}{|p{2.5cm}|p{2.5cm}|X|}
	\caption{Описание функций, используемых в модуле \texttt{macro.lsp}\label{macro:table}}\\
	\hline \centrow \setlength{\baselineskip}{0.7\baselineskip} Название функции & \centrow Аргументы функции & \centrow Описание функции \\
	\hline \centrow 1 & \centrow 2 & \centrow 3 \\ \hline
	\endfirsthead
	\caption*{Продолжение таблицы \ref{macro:table}}\\
	\hline \centrow 1 & \centrow 2 & \centrow 3 \\ \hline
	\finishhead
	add-fix-macro & name env arity args body & Сгенерированная макросом \texttt{mk/add-func} функция для сохранения информации о макросе с фиксированным числом аргументов \texttt{name} с окружением \texttt{env}, арностью \texttt{arity}, параметрами \texttt{args} и телом \texttt{body} в список \texttt{*fix-macros*}. \\ \hline
	add-nary-macro & name env arity args body & Сгенерированная макросом \texttt{mk/add-func} функция для сохранения информации о макросе с переменным числом аргументов \texttt{name} с окружением \texttt{env}, арностью \texttt{arity}, параметрами \texttt{args} и телом \texttt{body} в список \texttt{*nary-macros*}. \\ \hline
	make-macro & list & Сохранить информацию о макросе. \texttt{list} -- \texttt{(name args body)}, где \texttt{name} -- название макроса, \texttt{args} -- список параметров макроса, \texttt{body} -- тело макроса. \\ \hline
	subst & sym env & Подставить значение символа или глобальной переменной \texttt{sym} с окружением \texttt{env}. \\ \hline
	macro-eval & expr env & Рекурсивное раскрытие макроса \texttt{expr} с окружением \texttt{env}. \\ \hline
	macro-eval-backquote & expr env & Макровычисление квазицитирования \texttt{expr} с окружением \texttt{env}. \\ \hline
	macro-eval-if & expr env & Макровычисление условия \texttt{expr} с окружением \texttt{env}. \\ \hline
	macro-eval-args & args env & Макровычисление списка аргументов \texttt{args} с окружением \texttt{env}. \\ \hline
	macro-eval-progn & expr env & Макровычисление последовательности \texttt{expr} с окружением \texttt{env}. \\ \hline
	extend-macro-env & env args vals & Расширение окружения подстановки макросов \texttt{env} символами \texttt{args} и их значениями \texttt{vals}. \\ \hline
	macro-eval-let & args env & Макровычисление let-выражения \texttt{args} с временным расширением окружения \texttt{env}. \\ \hline
	macro-eval-app-func & args body vals env & Применение пользовательской функции с параметрами \texttt{args}, телом \texttt{body}, аргументами \texttt{vals} и окружением \texttt{env} внутри макроса. \\ \hline
	check-arguments & f type count args & Проверка на то, для функции или примитива \texttt{f} типа \texttt{type} список аргументов \texttt{args} имеет правильную длину в сравнении с числом фиксированных аргументов \texttt{count}. \\ \hline
	macroexpand & args vals body & Раскрытие макроса с параметрами \texttt{args}, аргументами \texttt{vals} и телом \texttt{body}. \\ \hline
	macro-eval-app & f args env & Макро применение функции, макроса или примитива \texttt{f} с аргументами \texttt{args} и окружением \texttt{env}. \\ \hline
	macro-eval-funcall & f args env & Макровычисление вызова функции \texttt{f} с аргументами \texttt{args} и окружением \texttt{env}. \\ \hline
	macro-eval-cond & list env & Макровычисление cond-выражения \texttt{list} с окружением \texttt{env}. \\ \hline
	macro-eval-setq & var expr env & Присвоение макро символу \texttt{var} значения \texttt{expr} с окружением \texttt{env}. \\ \hline
	macro-eval-function & s & Макровычисление функционального объекта \texttt{s}. \\ \hline
	macro-expand-progn & expr env & Макро раскрытие последовательности \texttt{expr} с окружением \texttt{env}. \\ \hline
\end{xltabular}
\renewcommand{\arraystretch}{1.0} % восстановление сетки


\subsection{compiler.lsp}

\texttt{compiler.lsp} -- модуль, отвечающий за обработку исходного S-выражения и его преобразование в промежуточное представление, состоящее из более простых операций. Собирает вспомогательную информацию, такую как число использования тех или иных функций, наличие рекурсии в функциях и т.д., для дальнейшей оптимизации.

Описание функций модуля \texttt{compiler.lsp} представлено в таблице \ref{compiler:table}.

\renewcommand{\arraystretch}{0.8} % уменьшение расстояний до сетки таблицы
\begin{xltabular}{\textwidth}{|p{2.5cm}|p{2.5cm}|X|}
	\caption{Описание функций, используемых в модуле \texttt{compiler.lsp}\label{compiler:table}}\\
	\hline \centrow \setlength{\baselineskip}{0.7\baselineskip} Название функции & \centrow Аргументы функции & \centrow Описание функции \\
	\hline \centrow 1 & \centrow 2 & \centrow 3 \\ \hline
	\endfirsthead
	\caption*{Продолжение таблицы \ref{compiler:table}}\\
	\hline \centrow 1 & \centrow 2 & \centrow 3 \\ \hline
	\finishhead
	extend-env & env args & Расширяет окружение \texttt{env} новым кадром, состоящим из аргументов \texttt{args}. \\ \hline
	add-func & name env arity args body & Сгенерированная макросом \texttt{mk/add-func} функция для сохранения информации о пользовательских функциях с фиксированным числом аргументов \texttt{name} с окружением \texttt{env}, арностью \texttt{arity}, параметрами \texttt{args} и телом \texttt{body} в список \texttt{*fix-functions*}. \\ \hline
	add-nary-func & name env arity args body & Сгенерированная макросом \texttt{mk/add-func} функция для сохранения информации о пользовательских функциях с переменным числом аргументов \texttt{name} с окружением \texttt{env}, арностью \texttt{arity}, параметрами \texttt{args} и телом \texttt{body} в список \texttt{*nary-functions*}. \\ \hline
	add-local-func & name env arity args body & Сгенерированная макросом \texttt{mk/add-func} функция для сохранения информации о локальных функциях \texttt{name} с окружением \texttt{env}, арностью \texttt{arity} и меткой \texttt{real-name} в список \texttt{*local-functions*}. \\ \hline
	inner-compile & expr env tail & Рекурсивная функция компиляции выражения \texttt{expr} c окружением \texttt{env} в промежуточную форму. \texttt{tail} -- является ли текущее выражение хвостовым. \\ \hline
	compile-func-body & name args body env & Компиляция тела функции \texttt{name} с параметрами \texttt{args}, телом \texttt{body} и окружением \texttt{env}. \\ \hline
	compile-lambda & name args body env & Регистрирует функцию \texttt{name} и компилирует её тело \texttt{body}. \texttt{args} -- параметры функции, \texttt{env} -- окружение. \\ \hline
	find-func & f &  Определение типа функции \texttt{f} и поиск её информации в соответствующем списке. \\ \hline
	compile-progn & lst env tail & Компилирует блок \texttt{progn} со списком выражений \texttt{lst} и окружением \texttt{env}. \texttt{tail} -- является ли текущее выражение хвостовым. \\ \hline
	inc-function-ref & name & Увеличить число ссылок на функцию \texttt{name} \\ \hline
	compile-application & f args env tail & Компиляция применения функции \texttt{f} с аргументами \texttt{args} в окружении \texttt{env}. \texttt{tail} -- является ли текущее выражение хвостовым. \\ \hline
	compile-function & f env & Компиляция создания замыкания функции \texttt{f} с окружением \texttt{env}. \\ \hline
	compile-defun & expr env & Компилирует объявление функции \texttt{expr} с помощью \texttt{defun} в окружении \texttt{env}. \\ \hline
	add-global & sym & Добавляет глобальную переменную \texttt{sym} и возвращает её индекс в списке глобальных переменных \texttt{*global-variables*}. \\ \hline
	find-global-var & var & Производит поиск переменной в глобальном окружении по символу \texttt{var}. Возвращает индекс в списке глобального окружения, если символ найден, иначе nil. \\ \hline
	var-search & params var & Вычисление порядкового номера переменной \texttt{var} в списке \texttt{params} с учетом \texttt{\&rest}. \\ \hline
	find-local-var & var env & Производит поиск переменной в локальном окружении \texttt{env} по символу \texttt{var}. Возвращает индекс переменной в окружении, если символ найден, иначе \texttt{nil}. \\ \hline
	find-var & var env & Производит поиск переменной по символу \texttt{var} сначала в локальном окружении \texttt{env}, а затем в списке глобальных переменных \texttt{*global-variables*}. Возвращает список, состоящий из символа типа переменной (\texttt{GLOBAL}, \texttt{LOCAL} или \texttt{DEEP}) и индекса переменной. \\ \hline
	compile-setq & body env & Компилирует setq-выражение с телом \texttt{body} и окружением \texttt{env}. \\ \hline
	compile-if & expr env tail & Компилирует if-выражение с телом \texttt{expr} и окружением \texttt{env}. \texttt{tail} -- является ли текущее выражение хвостовым. \\ \hline
	compile-constant & c & Компиляция константы \texttt{c}. \\ \hline
	compile-labels & lst env tail & Компилирует блок \texttt{labels} с телом \texttt{lst} и окружением \texttt{env}. \texttt{tail} -- является ли текущее выражение хвостовым. \\ \hline
	compile-variable & v env & Компиляция переменной \texttt{v} в окружении \texttt{env}. \\ \hline
	compile-backquote & expr env & Компиляция квазицитирования с телом \texttt{expr} и окружением \texttt{env}. \\ \hline
	compile-tagbody & body env & Компиляция \texttt{tagbody} с телом \texttt{body} и окружением \texttt{env}. \\ \hline
	compile-catch & args env & Компиляция \texttt{catch} с телом \texttt{args} и окружением \texttt{env}. \\ \hline
	compile-throw & args env & Компиляция \texttt{throw} с телом \texttt{args} и окружением \texttt{env}. \\ \hline
	compile-defconst & body env tail & Компиляция \texttt{defconst} с телом \texttt{body} и окружением \texttt{env} и сохранение этой константы в хеше \texttt{*const-vals*}. \texttt{tail} -- является ли текущее выражение хвостовым. \\ \hline
	compile-apply & func args pack env tail & Компиляция применения функционального объекта \texttt{func} с аргументами \texttt{args} и окружением \texttt{env}. \texttt{pack} -- нужно ли собирать аргументы в список (для примитива \texttt{apply}), \texttt{tail} -- является ли текущее выражение хвостовым. \\ \hline
	compile-args & args env & Компиляция аргументов \texttt{args} в окружении \texttt{env}. \\ \hline
	compile & expr & Анализ S-выражения \texttt{expr} и преобразование в эквивалентное выражение, состоящее из более простых операций.
\end{xltabular}
\renewcommand{\arraystretch}{1.0} % восстановление сетки


\subsection{optimizer.lsp}

\texttt{optimizer.lsp} -- модуль, отвечающий за оптимизацию промежуточной формы. Осуществляет обход дерева в глубину.

Описание функций модуля \texttt{optimizer.lsp} представлено в таблице \ref{optimizer:table}.

\renewcommand{\arraystretch}{0.8} % уменьшение расстояний до сетки таблицы
\begin{xltabular}{\textwidth}{|p{2.5cm}|p{2.5cm}|X|}
	\caption{Описание функций, используемых в модуле \texttt{optimizer.lsp}\label{optimizer:table}}\\
	\hline \centrow \setlength{\baselineskip}{0.7\baselineskip} Название функции & \centrow Аргументы функции & \centrow Описание функции \\
	\hline \centrow 1 & \centrow 2 & \centrow 3 \\ \hline
	\endfirsthead
	\caption*{Продолжение таблицы \ref{optimizer:table}}\\
	\hline \centrow 1 & \centrow 2 & \centrow 3 \\ \hline
	\finishhead
	accumulator-loading & tree & Удаляет избыточные формы загрузки значений в аккумулятор. \texttt{tree} -- форма \texttt{SEQ}. \\ \hline
	trivial-condition & alter & Если условие формы \texttt{ALTER} -- константа, то считает значение условия и сокращает форму \texttt{ALTER} до одной из её веток. \texttt{alter} -- форма \texttt{ALTER}. \\ \hline
	simplify-arithmetic & tree & Заменяет вызов арифметического примитива \texttt{tree} с неограниченным количеством аргументов на вложенные вызовы примитивов с фиксированным количеством аргументов. \texttt{tree} -- вызов примитива с переменным числом аргументов. \\ \hline
	constant-folding & tree & Свёртка и распространение констант. \texttt{tree} -- константа или вызов примитива. \\ \hline
	expand-seqs & tree & Разворачивает вложенные \texttt{SEQ}. \texttt{tree} -- форма \texttt{SEQ}. \\ \hline
	dead-code-elimination & tree & Убирает из формы \texttt{SEQ} неиспользуемый код (последовательность команд загрузки аккумулятора (кроме последней команды и команды перед \texttt{RETURN})). \\ \hline
	tail-call & tree & Оптимизация хвостовой рекурсии. \texttt{tree} -- форма \texttt{TAIL-CALL} или \texttt{TAIL-NCALL}. \\ \hline
	unused-functions & tree & Удаление неиспользуемых функций. \texttt{tree} -- форма \texttt{FUNC}. \\ \hline
	search-abs-tree & exp nodes & Возвращает истину если найдены хотя бы один узел из списка nodes. \\ \hline
	one-ref-args & body & Проверка на то, что в теле функции \texttt{body} к каждому аргументу только одно обращение. \\ \hline
	search-closure-deep-ref & exp & Поиск всех \texttt{FIX-CLOSURE} и проверка каждого тела замыкания на наличие \texttt{DEEP-REF}. \\ \hline
	beta-exp & exp args level & Выполнить подстановку тела функции \texttt{exp} с аргументами \texttt{args} и номером кадра активации \texttt{level}. \\ \hline
	make-nary-param & num args & Создать выражение для подстановки аргументов \texttt{args}, \texttt{num} -- число фиксированных аргументов. \\ \hline
	beta-expansion & call & Подстановка тела функции \texttt{call}, когда функция вызывается один раз и это возможно. \\ \hline
	can-inline & f & Возможно ли встраивание функции \texttt{f}. \\ \hline
	side-effects & exp & Есть ли в выражении \texttt{exp} побочные эффекты. \\ \hline
	optimize-tree1 & tree & Первый проход оптимизации промежуточной формы \texttt{tree}. Возвращает оптимизированное дерево. \\ \hline
	optimize-tree2 & tree & Второй проход оптимизатора -- встраивание функций, удаление после встраивания. \\ \hline
	optimize-tree & tree & Оптимизация промежуточной формы \texttt{tree}. \\ \hline
\end{xltabular}
\renewcommand{\arraystretch}{1.0} % восстановление сетки


\subsection{generator.lsp}

\texttt{generator.lsp} -- модуль, отвечающий за генерацию линейного кода ассемблера, состоящего из меток и инструкций с мнемониками и операндами.

Описание функций модуля \texttt{generator.lsp} представлено в таблице \ref{generator:table}.

\renewcommand{\arraystretch}{0.8} % уменьшение расстояний до сетки таблицы
\begin{xltabular}{\textwidth}{|p{2.5cm}|p{2.5cm}|X|}
	\caption{Описание функций, используемых в модуле \texttt{generator.lsp}\label{generator:table}}\\
	\hline \centrow \setlength{\baselineskip}{0.7\baselineskip} Название функции & \centrow Аргументы функции & \centrow Описание функции \\
	\hline \centrow 1 & \centrow 2 & \centrow 3 \\ \hline
	\endfirsthead
	\caption*{Продолжение таблицы \ref{generator:table}}\\
	\hline \centrow 1 & \centrow 2 & \centrow 3 \\ \hline
	\finishhead
	emit & val & Добавить инструкцию \texttt{val} в массив линейных инструкций \texttt{*program*}. \\ \hline
	inner-generate & expr & Рекурсивная генерация кода для скомпилированного выражения \texttt{expr}. \\ \hline
	generate-if & expr & Генерация ветвления. \texttt{expr} -- список из условия, ветки по истине и ветки по лжи. \\ \hline
	generate-2-params & expr & Генерация функции \texttt{expr} с 2-мя параметрами. \\ \hline
	generate-deep & i j expr & Генерация формы \texttt{DEEP-SET}. \texttt{i} -- смещение в окружении, \texttt{j} -- индекс в кадре активации, \texttt{expr} -- присваиваемое выражение. \\ \hline
	generate-set & expr & Генерация форм \texttt{GLOBAL-SET} и \texttt{LOCAL-SET}. \\ \hline
	generate-args & args set rev & Генерация вычисления аргументов \texttt{args}. \texttt{set} -- тип инструкции для каждого аргумента (например, \texttt{PUSH}), \texttt{rev} -- в обратном порядке. \\ \hline
	generate-fix-prim & type args & Генерация вызова примитива с фиксированным числом аргументов \texttt{type} и аргументами \texttt{args}. \\ \hline
	generate-nary-prim & type num args & Генерация вызова примитива с переменным числом аргументов \texttt{type} и аргументами \texttt{args}. \texttt{num} -- число фиксированных аргументов. \\ \hline
	generate-reg-call & name fix-num env args & Генерация обычного вызова функции \texttt{name} с аргументами \texttt{args} и окружением \texttt{env}. \texttt{fix-num} -- число фиксированных аргументов для функций с переменным числом аргументов. \\ \hline
	generate-let & count args body & Генерация let-формы с телом \texttt{body} и аргументами \texttt{args}, расширение окружения \texttt{env}. \\ \hline
	generate-closure & op name env body & Генерация замыкания. \texttt{op} -- инструкция создания замыкания (\texttt{PRIM-CLOSURE, \texttt{NPRIM-CLOSURE} или \texttt{FIX-CLOSURE}}), \texttt{name} -- название функционального объекта, env -- окружение, \texttt{body} -- опциональное тело (для лямбда-выражения). \\ \hline
	generate-catch & tag body & Генерация формы \texttt{catch} с меткой \texttt{tag} и телом \texttt{body}. \\ \hline
	generate-throw & tag res & Генерация формы \texttt{throw} с меткой \texttt{tag} и значением \texttt{res}. \\ \hline
	generate-tail-call & name fix-num depth args & Генерация хвостового вызова функции \texttt{name} с числом фиксированных аргументов \texttt{fix-num}, глубиной вызова относительно входа в функцию \texttt{depth} и аргументами \texttt{args}. \\ \hline
	generate-nary & args & Генерация списка необязательных аргументов \texttt{args}. \\ \hline
	generate-apply & func args pack & Генерация применения функции \texttt{func} к аргументам \texttt{args}. \texttt{pack} -- нужно ли объединять аргументы в список (для примитива \texttt{apply}). \\ \hline
	generate & expr & Генерация кода для скомпилированного выражения \texttt{expr}. \\ \hline
\end{xltabular}
\renewcommand{\arraystretch}{1.0} % восстановление сетки


\subsection{assembler.lsp}

\texttt{assembler.lsp} -- модуль, отвечающий за ассемблирование линейного кода в байт-код виртуальной машины, представляющий собой массив байт.

Описание функций модуля \texttt{assembler.lsp} представлено в таблице \ref{assembler:table}.

\renewcommand{\arraystretch}{0.8} % уменьшение расстояний до сетки таблицы
\begin{xltabular}{\textwidth}{|p{2.5cm}|p{2.5cm}|X|}
	\caption{Описание функций, используемых в модуле \texttt{assembler.lsp}\label{assembler:table}}\\
	\hline \centrow \setlength{\baselineskip}{0.7\baselineskip} Название функции & \centrow Аргументы функции & \centrow Описание функции \\
	\hline \centrow 1 & \centrow 2 & \centrow 3 \\ \hline
	\endfirsthead
	\caption*{Продолжение таблицы \ref{assembler:table}}\\
	\hline \centrow 1 & \centrow 2 & \centrow 3 \\ \hline
	\finishhead
	last-cdr & list & Найти последний \texttt{cdr} в списке \texttt{list}. \\ \hline
	asm-emit & \&rest bytes & Записать байты \texttt{bytes} в список \texttt{bytecode}. \\ \hline
	assemble-label & inst & Ассемблирование метки \texttt{inst}. \\ \hline
	assemble-const & inst & Ассемблирование инструкции загрузки константы \texttt{inst} в \texttt{ACC}. \\ \hline
	assemble-prim & inst & Ассемблирование инструкции вызова примитива \texttt{inst}. \\ \hline
	assemble-inst & inst jmp-labels & Ассемблирование инструкции общего вида \texttt{inst}. \\ \hline
	assemble-first-pass & program & Первый проход ассемблера (замена мнемоник инструкций соответствующими байтами и сбор информации о метках и переходах). \\ \hline
	assemble-second-pass & first-pass-res & Второй проход ассемблера (замена меток переходов на относительные адреса). \\ \hline
	assemble & program & Превращает список инструкций с метками \texttt{program} в байт-код. \\ \hline
\end{xltabular}
\renewcommand{\arraystretch}{1.0} % восстановление сетки


\subsection{Модульное тестирование фазы компиляции}

Модульные тесты фазы компиляции представлены в таблице \ref{unit_compiler:table}.

\renewcommand{\arraystretch}{0.8} % уменьшение расстояний до сетки таблицы
\begin{xltabular}{\textwidth}{|X|X|X|}
	\caption{Модульные тесты компилятора\label{unit_compiler:table}}\\
	\hline \centrow \setlength{\baselineskip}{0.7\baselineskip} Описание теста & \centrow \setlength{\baselineskip}{0.7\baselineskip} Исходное S-выражение & \centrow \setlength{\baselineskip}{0.7\baselineskip} Ожидаемый результат компиляции \\
	\hline \centrow 1 & \centrow 2 & \centrow 3 \\
	\endfirsthead
	\caption*{Продолжение таблицы \ref{unit_compiler:table}}\\
	\hline \centrow 1 & \centrow 2 & \centrow 3 \\ \hline
	\finishhead
	\hline Компиляция \texttt{progn} & \texttt{(progn 1 2 3)} & \texttt{(SEQ (CONST 1) (SEQ (CONST 2) (CONST 3)))} \\
	\hline Компиляция \texttt{if} & \texttt{(if t (progn 1 2) 3)} & \texttt{(ALTER (GLOBAL-REF 0) (SEQ (CONST 1) (CONST 2)) (CONST 3))} \\
	\hline Компиляция \texttt{setq} & \texttt{(setq a (* 2 3))} & \texttt{(GLOBAL-SET 2 (NARY-PRIM * 0 ((CONST 2) (CONST 3))))} \\
	\hline Компиляция объявления и вызова функции & \texttt{(progn (defun f (x) x) (f 5))} & \texttt{(SEQ (FUNC F (SEQ (LOCAL-REF 0) (RETURN))) (FIX-CALL F 0 ((CONST 5)) ()))} \\
	\hline Компиляция \texttt{let} и обращения к дальней переменной & \texttt{((lambda (x) ((lambda (y) (print x y)) 10)) 5)} & \texttt{(FIX-LET 1 ((CONST 5)) (FIX-LET 1 ((CONST 10)) (NARY-PRIM PRINT 0 ((DEEP-REF
		1 0) (LOCAL-REF 0)))))} \\
	\hline Компиляция вызова функционального объекта & \texttt{(funcall \#'+ 1 2)} & \texttt{(APPLY (NPRIM-CLOSURE +) ((CONST 1) (CONST 2)) T)} \\
	\hline Компиляция \texttt{tagbody} & \texttt{(tagbody (go end) 5 end)} & \texttt{(SEQ (GOTO END) (SEQ (CONST 5) (SEQ (LABEL END) (CONST NIL))))} \\
	\hline Компиляция \texttt{backquote} & \texttt{(progn (setq a 5) `(a = ,a))} & \texttt{(SEQ (GLOBAL-SET 2 (CONST 5)) (FIX-PRIM CONS ((CONST A) (FIX-PRIM CONS ((CONST =) (FIX-PRIM CONS ((GLOBAL-REF 2) (CONST ()))))))))} \\
	\hline Компиляция \texttt{catch} и \texttt{throw} & \texttt{(catch 'err (throw 'err "err"))} & \texttt{(CATCH (CONST ERR) (THROW (CONST ERR) (CONST "err")))} \\
\end{xltabular}
\renewcommand{\arraystretch}{1.0} % восстановление сетки


\subsection{Модульное тестирование фазы оптимизации}

Модульные тесты фазы оптимизации представлены в таблице \ref{unit_optimizer:table}.

\renewcommand{\arraystretch}{0.8} % уменьшение расстояний до сетки таблицы
\begin{xltabular}{\textwidth}{|X|X|X|}
	\caption{Модульные тесты оптимизатора\label{unit_optimizer:table}}\\
	\hline \centrow \setlength{\baselineskip}{0.7\baselineskip} Описание теста & \centrow \setlength{\baselineskip}{0.7\baselineskip} Исходное выражение & \centrow \setlength{\baselineskip}{0.7\baselineskip} Ожидаемое оптимизированное выражение \\
	\hline \centrow 1 & \centrow 2 & \centrow 3 \\
	\endfirsthead
	\caption*{Продолжение таблицы \ref{unit_optimizer:table}}\\
	\hline \centrow 1 & \centrow 2 & \centrow 3 \\ \hline
	\finishhead
	\hline Упрощение тривиального условия & \texttt{(ALTER (CONST T) (CONST 1) (CONST 2))} & \texttt{(CONST 1)} \\
	\hline Оптимизация арифметического выражения & \texttt{(NARY-PRIM + 0 ((LOCAL-REF 0) (LOCAL-REF 1)))} & \texttt{(FIX-PRIM + ((LOCAL-REF 0) (LOCAL-REF 1)))} \\
	\hline Оптимизация загрузки аккумулятора & \texttt{(SEQ (CONST 1) (CONST 2) (GLOBAL-SET 2 (CONST 3)) (CONST 4) (GLOBAL-REF 2))} & \texttt{(SEQ (GLOBAL-SET 2 (CONST 3)))} \\
	\hline Свёртка констант & \texttt{(GLOBAL-SET 2 (FIX-PRIM + ((CONST 1) (CONST 2))))} & \texttt{(GLOBAL-SET 2 (CONST 3))} \\
	\hline Встраивание и удаление функции & \texttt{(SEQ (FUNC LIST (SEQ (SEQ (CONST "Функция создания списка") (LOCAL-REF 0)) (RETURN))) (NARY-CALL LIST 0 0 ((CONST 1) (CONST 2) (CONST 3)) ()))} & \texttt{(NARY ((CONST 1) (CONST 2) (CONST 3)))} \\
	\hline Встраивание функции и удаление неиспользуемых функций & \texttt{(SEQ (FUNC F1 (SEQ (CONST 10) (RETURN))) (FUNC F2 (SEQ (CONST 10) (RETURN))) (FIX-CALL F2 0 () ()))} & \texttt{(CONST 10)} \\
	\hline Встраивание функций, удаление неиспользуемых функций и свёртка констант & \texttt{(SEQ (FUNC LIST (SEQ (SEQ (CONST "Функция создания списка") (LOCAL-REF 0)) (RETURN))) (SEQ (FUNC F (SEQ (NARY-PRIM + 0 ((NARY-PRIM * 0 ((FIX-PRIM CAR ((NARY-CALL LIST 0 0 ((CONST 1) (CONST 2) (CONST 3)) ()))) (CONST 2))) (LOCAL-REF 0) (CONST 10))) (RETURN))) (FIX-CALL F 0 ((CONST 5)) ())))} & \texttt{(CONST 17)} \\
\end{xltabular}
\renewcommand{\arraystretch}{1.0} % восстановление сетки


\subsection{Системное тестирование разработанного компилятора}

Было проведено системное тестирование разработанного компилятора, покрывающее все фазы компиляции и проверяющее корректность работы скомпилированного кода на тестах библиотек Common Lisp, исчерпывающе покрывающих основные граничные случаи компилятора. Успешным результатом данного тестирования является факт успешного выполнения тестов библиотек на виртуальной машине с помощью байт-кода, сгенерированного компилятором.


\subsubsection{Технология системного тестирования}

Для тестирования написана программа на языке C, реализующая виртуальную машину и читающая файл из стандартного потока ввода. Формат входного файла должен включать следующие элементы, разделённые пробелами:

\begin{itemize}
\item размер байт-кода N (суммарное количество опкодов и операндов);
\item непосредственно байт-код размером N (N чисел);
\item массив констант (например, \texttt{\#(T () 0 1 ...)});
\item размер массива глобальных переменных;
\item хеш-таблица меток и их адресов (например, \texttt{((ABS . 2) (G314 . 22) (NOT . 27) ...)}).
\end{itemize}

Пример файла скомпилированной программы \texttt{(print 5)}:

\begin{figure}[H]
\begin{lstlisting}[language={}]
8 0 2 10 11 1 19 9 20 #( T () 5 ) 2 ()
\end{lstlisting}
\end{figure}

Здесь \texttt{8} -- размер байт-кода, затем идёт сам байт-код, представляющий собой 8 чисел, потом массив констант -- \texttt{\#( T () 5 )}, размер массива глобальных переменных -- \texttt{2} и хеш-таблица меток и адресов -- \texttt{()}, в данном случае пустая. Таким образом, программа на C загружает все данные из файла, инициализирует состояние виртуальной машины и запускает выполнение байт-кода.

Для удобного автоматического тестирования скомпилирована версия самого компилятора, который загружает S-выражение из массива констант и выполняет дальнейшую компиляцию и печать бинарных данных в выше указанном формате. Константа исходного S-выражения в бинарном файле компилятора имеет следующий вид:

\begin{figure}[H]
\begin{lstlisting}[language=Lisp]
(PROGN *TARGET*)
\end{lstlisting}
\end{figure}

При подстановке S-выражения, которое требуется скомпилировать, вместо слова \texttt{*TARGET*}, и выполнения этого бинарного файла на виртуальной машине на выходе появляется бинарный файл скомпилированного исходного S-выражения. Если теперь этот файл выполнить на виртуальной машине, то произойдёт выполнение исходного S-выражения.

С помощью программных утилит, например \texttt{coreutils} (\texttt{cat}, \texttt{sed}), и перенаправления стандартных потоков ввода и вывода можно автоматизировать весь этот процесс компиляции произвольных программ на Common Lisp и их выполнения на виртуальной машине с помощью, например, скриптов \texttt{bash} или \texttt{Makefile}.


\subsubsection{Тестирование <<AES>>}

\texttt{aes} -- библиотека, реализующая алгоритм AES для шифрования и дешифрования данных. В ходе тестирования проверяются:

\begin{itemize}
\item процедура расширения ключа для AES-128;
\item шифрование и дешифрование AES-128 без дополнения;
\item шифрование и дешифрование с дополнением по схеме CMS;
\item корректность формирования и удаления CMS-дополнения при шифровании.
\end{itemize}

Результаты проведения тестирования на виртуальной машине представлены на рисунке \ref{test-aes:image}.

\begin{figure}[H]
\begin{lstlisting}[language=Lisp]
"Тестирование расширения ключа AES-128"
(OK #(#(0 1 2 3 4 5 6 7 8 9 10 11 12 13 14 15) #(214 170 116 253 210 175 114 250 218 166 120 241 214 171 118 254) #(182 146 207 11 100 61 189 241 190 155 197 0 104 48 179 254) #(182 255 116 78 210 194 201 191 108 89 12 191 4 105 191 65) #(71 247 247 188 149 53 62 3 249 108 50 188 253 5 141 253) ...) #(#(0 1 2 3 4 5 6 7 8 9 10 11 12 13 14 15) #(214 170 116 253 210 175 114 250 218 166 120 241 214 171 118 254) #(182 146 207 11 100 61 189 241 190 155 197 0 104 48 179 254) #(182 255 116 78 210 194 201 191 108 89 12 191 4 105 191 65) #(71 247 247 188 149 53 62 3 249 108 50 188 253 5 141 253) ...))
"Тесты для AES-128 шифрования и дешифрования"
(OK #(15 14 13 12 11 10 9 8 7 6 5 4 3 2 1 0) #(15 14 13 12 11 10 9 8 7 6 5 4 3 2 1 0))
(OK #(11 14 12 9 1 6 8 13 3 15 2 0 10 7 5 4) #(11 14 12 9 1 6 8 13 3 15 2 0 10 7 5 4))
"Тестирование шифрования и дешифрования данных с паддингом CMS"
(OK #(1 2 3 4 5 6 7) #(1 2 3 4 5 6 7))
(OK #(0 1 2 3 4 5 6 7 8 9 10 11 12 13 14 15 16 17 18 19) #(0 1 2 3 4 5 6 7 8 9 10 11 12 13 14 15 16 17 18 19))
"Тест правильности паддинга CMS при шифровании данных"
(OK #(10 20 30 40 50) #(10 20 30 40 50))
"Успешно завершено тестов: " 6 / 6
"Все тесты завершены успешно"
\end{lstlisting}
\caption{Результаты выполнения тестов AES на виртуальной машине}
\label{test-aes:image}
\end{figure}


\subsubsection{Тестирование <<regex>>}

\texttt{regex} -- библиотека регулярных выражений. В ходе тестирования проверяется парсинг regex-шаблонов и множество случаев поиска regex-шаблона в строке.

Результаты проведения тестирования на виртуальной машине представлены на рисунке \ref{test-regex:image}.

\begin{figure}[H]
\begin{lstlisting}[language=Lisp]
"Тестирование символа"
(OK (SYM #\a) (SYM #\a))
"Тестирование количества"
(OK (REPEAT 12 ()) (REPEAT 12 ()))
(OK (REPEAT 12 MORE) (REPEAT 12 MORE))
(OK (REPEAT 12 50) (REPEAT 12 50))
"Тестирование количества"
(OK (RANGE () ((65 ()) (CLASS DIGIT) (CLASS NSPACE) (45 ()) (98 ()))) (RANGE () ((65 ()) (CLASS DIGIT) (CLASS NSPACE) (45 ()) (98 ()))))
(OK (RANGE NEGATIVE ((65 ()))) (RANGE NEGATIVE ((65 ()))))
(OK (RANGE () ((65 122) (66 115))) (RANGE () ((65 122) (66 115))))
"Тестирование квантификатора"
(OK ((REPEAT 12 ()) LAZY) ((REPEAT 12 ()) LAZY))
(OK ((PLUS) LAZY) ((PLUS) LAZY))
(OK ((STAR) LAZY) ((STAR) LAZY))
"Тестирование элемента"
(OK (LAZY-STAR (SYM #\+)) (LAZY-STAR (SYM #\+)))
(OK (OR (EPSILON) (SYM #\b)) (OR (EPSILON) (SYM #\b)))
(OK (SEQ (SYM #\b) (STAR (...))) (SEQ (SYM #\b) (STAR (SYM #\b))))
(OK (LAZY-STAR (ANY)) (LAZY-STAR (ANY)))
...
"Тестирование регулярных выражений"
(OK "" "") STRING "a" REGEX "a{0}" GROUPS ()
(OK "abccccd" "abccccd") STRING "abccccd" REGEX "abc{4}d" GROUPS ()
(OK "assaaabbaa" "assaaabbaa") STRING "assaaabbaa" REGEX "a+(s*)(a+(b+)a+)" GROUPS ("ss" "aaabbaa" "bb")
...
(OK "привет" "привет") STRING "abc_123привет" REGEX "[^A-Za-z0-9_]+" GROUPS ()
(OK ".*?" ".*?") STRING ".*?" REGEX "\.\*\?" GROUPS ()
(OK "nonnoncapturinggroups" "nonnoncapturinggroups") STRING "nonnoncapturinggroups" REGEX "(?:non)+(?:capturing)(?:groups)" GROUPS ()
(OK "134" "134") STRING "_134_" REGEX "(1)(2|34|9)" GROUPS ("1" "34")
(OK "1345" "1345") STRING "_1345_" REGEX "(1)(34|3|2)(45)" GROUPS ("1" "3" "45")
...
"Успешно завершено тестов: " 76 / 76
"Все тесты завершены успешно"
\end{lstlisting}
\caption{Результаты выполнения тестов regex на виртуальной машине}
\label{test-regex:image}
\end{figure}


\subsubsection{Тестирование <<PNG>>}

\texttt{png} -- библиотека для парсинга файлов изображений в формате PNG. В ходе тестирования проверяется успешный разбор PNG-файла.

Результаты проведения тестирования на виртуальной машине представлены на рисунке \ref{test-png:image}.

\begin{figure}[H]
\begin{lstlisting}[language=Lisp]
(OK #(#(#(145 99 170 255) #(119 85 177 255) #(189 75 108 255) #(251 186 181 255) #(69 117 255 255) #(176 189 168 255) #(255 235 59 255) #(255 235 59 255)) #(#(41 98 255 255) #(118 117 163 255) #(42 98 254 255) #(97 165 154 255) #(64 148 146 255) #(193 218 108 255) #(255 237 74 255) #(255 237 78 255)) #(#(41 98 255 255) #(229 146 31 255) #(73 105 217 255) #(42 99 253 255) #(76 175 80 255) #(103 187 106 255) #(150 178 255 255) #(66 115 255 255)) #(#(90 134 255 255) #(235 150 32 255) #(99 129 223 255) #(100 162 254 255) #(128 216 255 255) #(128 216 255 255) #(119 205 255 255) #(96 172 255 255)) #(#(131 217 255 255) #(241 159 28 255) #(158 201 195 255) #(137 218 240 255) #(160 221 206 255) #(160 221 206 255) #(160 221 206 255) #(160 221 206 255)) #(#(149 222 255 255) #(241 161 32 255) #(178 218 210 255) #(254 236 64 255) #(255 235 59 255) #(255 235 59 255) #(255 235 59 255) #(255 235 59 255))) #(#(#(145 99 170 255) #(119 85 177 255) #(189 75 108 255) #(251 186 181 255) #(69 117 255 255) #(176 189 168 255) #(255 235 59 255) #(255 235 59 255)) #(#(41 98 255 255) #(118 117 163 255) #(42 98 254 255) #(97 165 154 255) #(64 148 146 255) #(193 218 108 255) #(255 237 74 255) #(255 237 78 255)) #(#(41 98 255 255) #(229 146 31 255) #(73 105 217 255) #(42 99 253 255) #(76 175 80 255) #(103 187 106 255) #(150 178 255 255) #(66 115 255 255)) #(#(90 134 255 255) #(235 150 32 255) #(99 129 223 255) #(100 162 254 255) #(128 216 255 255) #(128 216 255 255) #(119 205 255 255) #(96 172 255 255)) #(#(131 217 255 255) #(241 159 28 255) #(158 201 195 255) #(137 218 240 255) #(160 221 206 255) #(160 221 206 255) #(160 221 206 255) #(160 221 206 255)) #(#(149 222 255 255) #(241 161 32 255) #(178 218 210 255) #(254 236 64 255) #(255 235 59 255) #(255 235 59 255) #(255 235 59 255) #(255 235 59 255))))
()
"Успешно завершено тестов: " 1 / 1
"Все тесты завершены успешно"
\end{lstlisting}
\caption{Результаты выполнения тестов PNG на виртуальной машине}
\label{test-png:image}
\end{figure}


\subsubsection{Тестирование <<JPEG>>}

\texttt{jpeg} -- библиотека для парсинга файлов изображений в формате JPEG. В ходе тестирования проверяется успешный разбор JPEG-файла.

Результаты проведения тестирования на виртуальной машине представлены на рисунке \ref{test-jpeg:image}.

\begin{figure}[H]
\begin{lstlisting}[language=Lisp]
P3 113 115 255
#(#(0 87 36 0 87 36 0 87 36 0 87 36 0 87 36 0 87 36 0 87 36 0 87 36 0 87 36 0 87 36 0 87 36 0 87 36 0 87 36 0 87 36 0 87 36 0 87 36 0 87 36 0 87 36 0 87 36 0 87 36 0 87 36 0 87 36 0 87 36 0 87 36 0 87 36 0 87 36 0 87 36 0 87 36 0 87 36 0 87 36 0 87 36 0 87 36 0 87 36 0 87 36 0 87 36 0 87 36 0 87 36 0 87 36 0 87 36 0 87 36 0 87 36 0 87 36 0 87 36 0 87 36 0 87 36 0 87 36 0 87 36 0 87 36 0 90 36 0 89 36 0 88 36 0 87 36 3 85 36 6 84 36 10 82 36 10 82 36 13 80 36 12 81 36 10 82 36 9 82 36 7 83 36 6 84 36 5 84 36 3 85 36 0 87 36 0 87 36 0 87 36 0 87 36 0 87 36 0 87 36 0 87 36 0 87 36 0 87 36 0 87 36 0 87 36 0 87 36 0 87 36 0 87 36 0 87 36 0 87 36 0 87 36 0 87 36 0 87 36 0 87 36 0 87 36 0 87 36 0 87 36 0 87 36 0 87 36 0 87 36 0 87 36 0 87 36 0 87 36 0 87 36 0 87 36 0 87 36 0 87 36 0 87 36 0 87 36 0 87 36 0 87 36 0 87 36 0 87 36 0 87 36 0 87 36 0 87 36 0 87 36 0 87 36 0 87 36 0 87 36 0 87 36 0 87 36 0 87 36) #(0 87 36 0 87 36 0 87 36 0 87 36 0 87 36 0 87 36 0 87 36 0 87 36 0 87 36 0 87 36 0 87 36 0 87 36 0 87 36 0 87 36 0 87 36 0 87 36 0 87 36 0 87 36 0 87 36 0 87 36 0 87 36 0 87 36 0 87 36 0 87 36 0 87 36 0 87 36 0 87 36 0 87 36 0 87 36 0 87 36 0 87 36 0 87 36 0 87 36 0 87 36 0 87 36 0 87 36 0 87 36 0 87 36 0 87 36 0 87 36 0 87 36 0 87 36 0 87 36 0 87 36 0 87 36 0 87 36 0 87 36 0 87 36 0 89 36 0 89 36 0 88 36 0 87 36 3 85 36 6 84 36 9 82 36 9 82 36 12 81 36 10 82 36 9 82 36 7 83 36 6 84 36 5 84 36 3 85 36 2 86 36 0 87 36 0 87 36 0 87 36 0 87 36 0 87 36 0 87 36 0 87 36 0 87 36 0 87 36 0 87 36 0 87 36 0 87 36 0 87 36 0 87 36 ...
\end{lstlisting}
\caption{Результаты выполнения тестов JPEG на виртуальной машине}
\label{test-jpeg:image}
\end{figure}


\subsubsection{Тестирование <<Lua>>}

\texttt{lua} -- библиотека для трансляции Lua-скриптов в Common Lisp. В ходе тестирования проверяется лексический, синтаксический анализ и трансляция программы на языке Lua в аналогичную программу на Common Lisp.

Результаты проведения тестирования на виртуальной машине представлены на рисунке \ref{test-lua:image}.

\begin{figure}[H]
\begin{lstlisting}[language=Lisp]
"Выполняет тест парсинга лексемы переменной"
(OK ABCDF ABCDF)
"Выполняет тест парсинга лексемы числа в шестнадцетиричном виде"
(OK 175 175)
"Выполняет тест парсинга лексемы числа в шестнадцетиричном виде"
(OK 175 175)
"hash test"
2 2
(HASH ("a" "first") (1 "second") (2 "third") ("b" "fourth"))
"Выполняет тест парса функции"
(FUNCTION (LAMBDA (ARG) (PROGN (LUA-SET (ARG) (100)))))
"Выполняет тест лексического на Lua в Lisp-выражение"
(PRINT #\( A #\. B #\. C #\) A LUA-SET #\{ B LUA-SET "1" #\, C LUA-SET "2" #\, 3 #\, 4 #\} #\; ABC LUA-SET 100 END B LUA-SET "aboaasd" C LUA-SET FALSE-CONST LUA-AND NIL-CONST)
"Выполняет тест преобразования строки на Lua в Lisp-выражение"
(PROGN (LUA-SET (PERSON) ((LUA-CREATETABLE ()))) (LUA-SET-INDEX PERSON "NEW" (FUNCTION (LAMBDA (NAME AGE) (PROGN (LET ((P (LUA-CREATETABLE ()))) (PROGN (LUA-SET ((LUA-GET-INDEX P "NAME")) (NAME)) (LUA-SET ((LUA-GET-INDEX P "AGE")) (AGE)) (LUA-SET-INDEX P "GET_NAME" (FUNCTION (LAMBDA (SELF) (PROGN (LUA-GET-INDEX SELF "NAME"))))) (LUA-SET-INDEX P "SET_NAME" (FUNCTION (LAMBDA (SELF NAME) (PROGN (LUA-SET ((LUA-GET-INDEX SELF "NAME")) (NAME)))))) (LUA-SET-INDEX P "HELLO" (FUNCTION (LAMBDA (SELF) (PROGN (LUA-CONCAT (LUA-CONCAT (LUA-CONCAT (LUA-CONCAT "Hello, my name is " (LUA-GET-INDEX SELF "NAME")) ", i am ") (LUA-GET-INDEX SELF "AGE")) " years old"))))) P)))))) (LUA-SET (VASYA) ((FUNCALL (LUA-GET-INDEX PERSON "NEW") "Vasya" 30))) (LUA-SET (KOLYA) ((FUNCALL (LUA-GET-INDEX PERSON "NEW") "Kolya" 20))) (LUA-SET (VASYA KOLYA) (KOLYA VASYA)) (FUNCALL PRINT (LET ((__T VASYA)) (FUNCALL (LUA-GET-INDEX __T "HELLO") __T))) (FUNCALL PRINT (LET ((__T KOLYA)) (FUNCALL (LUA-GET-INDEX __T "HELLO") __T))))
"Выполняет тест исполнения Lisp-выражения, преобразованного из Lua кода"
(PROGN (LUA-SET (PERSON) ((LUA-CREATETABLE ()))) (LUA-SET-INDEX PERSON "NEW" (FUNCTION (LAMBDA (NAME AGE) (PROGN (LET ((P (LUA-CREATETABLE ()))) (PROGN (LUA-SET ((LUA-GET-INDEX P "NAME")) (NAME)) (LUA-SET ((LUA-GET-INDEX P "AGE")) (AGE)) (LUA-SET-INDEX P "GET_NAME" (FUNCTION (LAMBDA (SELF) (PROGN (LUA-GET-INDEX SELF "NAME"))))) (LUA-SET-INDEX P "SET_NAME" (FUNCTION (LAMBDA (SELF NAME) (PROGN (LUA-SET ((LUA-GET-INDEX SELF "NAME")) (NAME)))))) (LUA-SET-INDEX P "HELLO" (FUNCTION (LAMBDA (SELF) (PROGN (LUA-CONCAT (LUA-CONCAT (LUA-CONCAT (LUA-CONCAT "Hello, my name is " (LUA-GET-INDEX SELF "NAME")) ", i am ") (LUA-GET-INDEX SELF "AGE")) " years old"))))) P)))))) (LUA-SET (VASYA) ((FUNCALL (LUA-GET-INDEX PERSON "NEW") "Vasya" 30))) (LUA-SET (KOLYA) ((FUNCALL (LUA-GET-INDEX PERSON "NEW") "Kolya" 20))) (LUA-SET (VASYA KOLYA) (KOLYA VASYA)) (FUNCALL PRINT (LET ((__T VASYA)) (FUNCALL (LUA-GET-INDEX __T "HELLO") __T))) (FUNCALL PRINT (LET ((__T KOLYA)) (FUNCALL (LUA-GET-INDEX __T "HELLO") __T))))
"Успешно завершено тестов: " 3 / 3
"Все тесты завершены успешно"
\end{lstlisting}
\caption{Результаты выполнения тестов Lua на виртуальной машине}
\label{test-lua:image}
\end{figure}